%%%%%%%%%%  scan idx 115  |  folio 104  %%%%%%%%%%

Consideremos ahora la localización y la completación en la teoría de
deformaciones de esquemas. Fijemos primero la notación. Sea $X_0$ un
$k$-esquema de tipo finito y sea $x_0\in X_0$ un punto fijo. Existen
funtores naturales
\[
\underline M_{X_0}
\longrightarrow
\underline M_{\operatorname{Spec}(\mathcal O_{X_0,x_0})}
\longrightarrow
\underline M_{\operatorname{Spec}(\widehat{\mathcal O}_{X_0,x_0})},
\]

%%%%%%%%%%  scan idx 116  |  folio 105  %%%%%%%%%%

y, por tanto, funtores
\[
\widehat{\underline M}_{X_0}
\longrightarrow
\widehat{\underline M}_{\operatorname{Spec}(\mathcal O_{X_0,x_0})}
\longrightarrow
\widehat{\underline M}_{\operatorname{Spec}
 (\widehat{\mathcal O}_{X_0,x_0})}.
\]
Como se observó en 4.1(c), $\widehat{\underline M}_{X_0}$ es
canónicamente equivalente a la categoría cuyos objetos son pares
$(R,\mathfrak X)$, donde
$R\in\operatorname{Ob}(\widehat{\underline C}_\Lambda)$ y
$\mathfrak X$ es un esquema formal $R$-ádico plano provisto de un
isomorfismo fijo $\mathfrak X_0\xrightarrow{\sim}X_0$.

Cuando no haya peligro de confusión, denotaremos las imágenes de
$(R,\mathfrak X)$ por la localización y la completación mediante
$(R,\mathfrak X_{(x_0)})$ y
$(R,\mathfrak X_{(\widehat{x}_0)})$, respectivamente. Así, si
\[
\mathfrak X=\varinjlim_\nu\mathfrak X_\nu,
\qquad
\mathfrak X_\nu
=\mathfrak X\times_{\operatorname{Spf}(R)}
 \operatorname{Spec}(R/\mathfrak m^{\nu+1}),
\]
entonces
\[
\mathfrak X_{(x_0)}
=\varinjlim_\nu
 \operatorname{Spec}(\mathcal O_{\mathfrak X_\nu,x_0})
\]
y
\[
\mathfrak X_{(\sisi{\widehat{X}_0}{\widehat{x}_0})}
=\varinjlim_\nu
 \operatorname{Spec}(\widehat{\mathcal O}_{\mathfrak X_\nu,x_0}).
\]
